%%%%%
%
\section{Conversion d'un {\tt float HP} en {\tt float}}

\itemstyle{Op\'erateur~:} {\sf CnvHpI3e}
\index{CnvHpI3e}

%%%%%OBJET%%%%
\itemstyle{Objet~:}

Conversion d'un param\`etre {\tt float\_HP} en {\tt float} selon
la norme IEEE-754.

%%%%%ARGUMENTS%%%%
\bigskip
%%%%%ENTREES%%%%
\begin{tabarg}
\multicolumn{3}{l}{\titretab {Entr\'ees}}\\
\hline
$X_{hp}$ & Vecteur, \texttt{HPfloat32 } & Param\`etre \`a convertir (unit\'e indiff\'erente)\\
\end{tabarg}

\medskip
%%%%%SORTIES%%%%
\begin{tabarg}
\multicolumn{3}{l}{\titretab {Sorties}}\\
\hline
$X_{ieee}$ & Vecteur, \texttt{float} & Param\`etre transform\'e (m\^eme unit\'e
que $X_{hp}$)\\
\end{tabarg}


%%%%%FORMULE%%%%
\itemstyle{Formule}

%%%CnvHpI3e
\[
Resultat = a-b. \Zeta
\]



%%%%%AUTEUR%%%%
\itemstyle{Auteur~:}
C. Allet. 

%%%%%REFERENCE%%%%
\itemstyle{R\'ef\'erences~:}
Norme IEEE754 \cite{IEEE754}  \\ 
HP-UX Floating-Point Guide \cite{DOC-HPFLOAT},
Manuel utilisateur Turbo-C \cite{TURBO-C}, page 329.
